\section*{Bab \ref{ch:b1:field-lab-methods-1} --- Metode lapangan dan laboratorium I}

\begin{solution}{exo:b1:field-lab-methods-1:1}
Pada 1:10\,000, \qty{4.9}{m} menjadi $4.9/10\,000 = \qty{0.49}{mm}$; pada 1:25\,000, $4.9/25\,000 = \qty{0.20}{mm}$, kira-kira lebar satu garis pensil.
\end{solution}

\begin{solution}{exo:b1:field-lab-methods-1:2}
Di Reykjavik $D = \ang{-11.28}$: $52 - 11.28 = \ang{040.7}$. Di Wellington
$D = \ang{+23.23}$: $52 + 23.23 = \ang{075.2}$. Pembacaan kompas yang sama berarti arah yang berjarak \ang{34.5}.
\end{solution}

\begin{solution}{exo:b1:field-lab-methods-1:3}
$20\to015$. Dengan $t = \ang{15}$ dan $p = \ang{20}$, vektor satuan ke bawahnya adalah
\[
  (\cos\ang{20}\sin\ang{15},\ \cos\ang{20}\cos\ang{15},\ -\sin\ang{20}) =
  (0.2432,\ 0.9077,\ -0.3420).
\]
\end{solution}

\begin{solution}{exo:b1:field-lab-methods-1:4}
Rata-rata aritmetiknya adalah $(350 + 358 + 6)/3 = \ang{238}$, hampir berlawanan dengan ketiga pembacaan. Vektor satuan horizontalnya $(\sin a, \cos a)$ dijumlahkan menjadi $(-0.1736 - 0.0349 + 0.1045,\ 0.9848 + 0.9994 + 0.9945) = (-0.1040, 2.9787)$, dan \omterm{def:b1:spherical-astronomy:horizontal}{azimutnya} adalah $\operatorname{atan2}(-0.1040, 2.9787) = \ang{-2.0}$, yaitu \ang{358.0}: rata-rata yang benar, di tengah pembacaan.
\end{solution}

\begin{solution}{exo:b1:field-lab-methods-1:5}
Rata-rata $175/5 = \ang{35.0}$. Simpangan $-1, 2, 0, -2, 1$; jumlah kuadratnya 10;
$s = \sqrt{10/4} = \ang{1.58}$. Galat baku $1.58/\sqrt5 = \ang{0.71}$.
\end{solution}

\begin{solution}{exo:b1:field-lab-methods-1:6}
Normalnya adalah $(\sin\ang{30}, 0, \cos\ang{30}) = (0.5, 0, 0.866)$ dan
$(-0.5, 0, 0.866)$; $\vec S = (0, 0, 1.732)$ vertikal: rata-ratanya adalah bidang horizontal, dan $\bar R = 1.732/2 = 0.866$, dengan setiap pembacaan berjarak \ang{30} dari rata-rata. Kedua pembacaan tidak mencuplik satu permukaan dengan galat: keduanya termasuk dua sayap \omterm{def:b1:geological-maps:fold}{lipatan}, dua populasi yang berbeda, dan rata-ratanya tidak menggambarkan keduanya. Setiap sayap dirata-ratakan sendiri; bersama-sama mereka menentukan \omterm{def:b1:geological-maps:fold}{sumbu lipatan}, garis yang sama bagi kedua bidang, di sini horizontal dan utara--selatan.
\end{solution}

\begin{solution}{exo:b1:field-lab-methods-1:7}
$\tan p = \tan\ang{12}\cos(\ang{80} - \ang{120}) = 0.2126\times0.7660 =
0.1628$: $p = \ang{9.2}$, goresan itu menunjam ke arah \ang{120}. Untuk \omterm{def:b1:field-lab-methods-1:trend-plunge}{trend} \ang{170}, $\cos(\ang{-90}) = 0$: garisnya horizontal, sepanjang \omterm{def:b1:geological-maps:attitude}{jurus}.
\end{solution}

\begin{solution}{exo:b1:field-lab-methods-1:8}
$0.47/0.20 = 2.35$ tahun setelah 1 Januari 2026, yaitu musim semi 2028. Pada 1 Januari 2030: $-0.47 + 4\times0.20 = \ang{+0.33}$, ke timur dari utara geografis.
\end{solution}

\begin{solution}{exo:b1:field-lab-methods-1:9}
Ketiga normal terletak pada bidang vertikal berazimut \ang{080}, pada \ang{10}, \ang{12} dan \ang{14} dari vertikal; dua yang terluar simetris terhadap yang di tengah, sehingga jumlahnya berada di sepanjangnya: rata-ratanya adalah $080/12$, dengan $\theta = \ang{2}, 0, \ang{2}$. Maka $\bar R = (1 +
2\cos\ang{2})/3 = (1 + 1.998782)/3 = 0.999594$, dan $\sqrt{2(1 - \bar R)} =
\sqrt{0.000812} = \qty{0.02850}{rad} = \ang{1.633}$; secara eksak,
$\theta_{\text{rms}} = \sqrt{(4 + 0 + 4)/3} = \ang{1.633}$.
\end{solution}

\begin{solution}{exo:b1:field-lab-methods-1:10}
Penampang berupa lingkaran berjari-jari $\sqrt{r^2 - h^2}$. Rata-ratanya adalah
$\frac1r\int_0^r\sqrt{r^2 - h^2}\,\mathrm dh = \frac1r\cdot\frac{\pi r^2}{4}
= \frac{\pi r}{4}$, karena integralnya adalah luas seperempat cakram. Butir selebar \qty{0.50}{mm} rata-rata tampak selebar $0.50\times\pi/4 = \qty{0.39}{mm}$ dalam sayatan: \omterm{def:b1:field-lab-methods-1:thin-section}{sayatan tipis} meremehkan ukuran butir.
\end{solution}

\begin{solution}{exo:b1:field-lab-methods-1:11}
Rotasi sebesar $D$ terhadap vertikal,
$(x, y, z) \mapsto (x\cos D + y\sin D,\ -x\sin D + y\cos D,\ z)$, membawa
$(\sin\delta\sin\alpha, \sin\delta\cos\alpha, \cos\delta)$ menjadi
$(\sin\delta\sin(\alpha + D), \sin\delta\cos(\alpha + D), \cos\delta)$: ia menambahkan $D$ pada setiap \omterm{def:b1:geological-maps:attitude}{arah kemiringan}. Karena linear, ia membawa resultan $\vec S$ ke resultan normal-normal yang diputar; ia mempertahankan panjang dan komponen vertikal, sehingga \omterm{def:b1:geological-maps:attitude}{kemiringan} rata-rata $\arccos(S_U/\|\vec S\|)$ tidak berubah, dan ia memutar bagian horizontal $\vec S$ sebesar $D$, sehingga \omterm{def:b1:geological-maps:attitude}{arah kemiringan} rata-rata bertambah $D$. Mengoreksi setiap pembacaan terhadap deklinasi, atau hanya rata-ratanya, memberikan hasil yang sama.
\end{solution}

\begin{solution}{exo:b1:field-lab-methods-1:12}
Sebaran sepanjang garis \omterm{def:b1:geological-maps:attitude}{kemiringan} paling curam adalah sebaran \omterm{def:b1:geological-maps:attitude}{kemiringan}: galat bakunya $2/\sqrt N \le 0.5$ untuk $N \ge 16$. Melintangnya, putaran normal sebesar \ang{2} mengubah \omterm{def:b1:geological-maps:attitude}{arah kemiringan} sebesar $2/\sin\ang{12} = \ang{9.62}$; galat baku \omterm{def:b1:geological-maps:attitude}{arah kemiringan} rata-rata adalah $9.62/\sqrt N \le 1$ untuk $N \ge 92.5$, yaitu 93 pembacaan. \omterm{def:b1:geological-maps:attitude}{Kemiringan} yang landai cepat diketahui; arahnya tidak.
\end{solution}

\begin{solution}{pb:b1:field-lab-methods-1:1}
\textbf{1.} \omterm{def:b1:geological-maps:attitude}{Arah kemiringan}: $967/12 = \ang{80.6}$; \omterm{def:b1:geological-maps:attitude}{kemiringan}: $145/12 =
\ang{12.1}$.

\textbf{2.} \ang{355} dan \ang{005} rata-rata menjadi \ang{180}, arah yang berlawanan. Di sini semua \omterm{def:b1:geological-maps:attitude}{arah kemiringan} terletak antara \ang{076} dan \ang{085}, jauh dari lompatan di \ang{000}, sehingga rata-rata aritmetik berfungsi.

\textbf{3.} $\|\vec n\|^2 = \sin^2\delta(\sin^2\alpha + \cos^2\alpha) +
\cos^2\delta = 1$. Garis \omterm{def:b1:geological-maps:attitude}{kemiringan} ke bawah memiliki arah
$(\cos\delta\sin\alpha, \cos\delta\cos\alpha, -\sin\delta)$, dan hasil kali titiknya dengan $\vec n$ adalah $\sin\delta\cos\delta - \cos\delta\sin\delta = 0$; garis \omterm{def:b1:geological-maps:attitude}{jurus} adalah horizontal dan tegak lurus \omterm{def:b1:geological-maps:attitude}{arah kemiringan},
$(\cos\alpha, -\sin\alpha, 0)$, dan hasil kali titiknya dengan $\vec n$ adalah
$\sin\delta\sin\alpha\cos\alpha - \sin\delta\cos\alpha\sin\alpha = 0$.

\textbf{4.} Pembacaan 11 sama dengan pembacaan 2: $(0.2222, 0.0352, 0.9744)$.
Pembacaan 12:
\[
  (\sin\ang{11}\sin\ang{82},\ \sin\ang{11}\cos\ang{82},\ \cos\ang{11}) =
  (0.1890,\ 0.0266,\ 0.9816).
\]

\textbf{5.} $\vec S = (2.4747,\ 0.4082,\ 11.7294)$.

\textbf{6.} $\|\vec S\| = \sqrt{6.1241 + 0.1666 + 137.5788} = 11.99457$, dan
$\bar R = 11.99457/12 = 0.999547$.

\textbf{7.} $\bar\alpha = \operatorname{atan2}(2.4747, 0.4082) = \ang{80.6}$
dan $\bar\delta = \arccos(11.7294/11.99457) = \arccos 0.977893 = \ang{12.1}$:
rata-rata vektornya adalah $080.6/12.1$, sama dengan rata-rata aritmetiknya di sini.

\textbf{8.} \omterm{def:b1:geological-maps:attitude}{Arah kemiringan}: $s = \ang{2.94}$; \omterm{def:b1:geological-maps:attitude}{kemiringan}: $s = \ang{1.68}$.

\textbf{9.} $\theta_{\text{rms}} \approx \sqrt{2\times0.000453} =
\qty{0.0301}{rad} = \ang{1.72}$.

\textbf{10.} $\sqrt{1.68^2 + (\sin\ang{12.1}\times2.94)^2} = \sqrt{2.82 +
0.38} = \ang{1.79}$, dekat dengan \ang{1.72}: pembacaan berperilaku seperti yang dikatakan model galat, dalam batas efek pembulatan ke derajat dan sampel kecil.

\textbf{11.} \omterm{def:b1:geological-maps:attitude}{Kemiringan}: ia menyumbang 2.82 dari 3.20 (dalam derajat kuadrat). \omterm{def:b1:geological-maps:attitude}{Arah kemiringan} hanya berpengaruh melalui $\sin\delta = 0.21$, karena pada perlapisan yang landai perubahan \omterm{def:b1:geological-maps:attitude}{arah kemiringan} memutar bidang sangat sedikit.

\textbf{12.} \omterm{def:b1:geological-maps:attitude}{Kemiringan} rata-rata: $1.68/\sqrt{12} = \ang{0.48}$; \omterm{def:b1:geological-maps:attitude}{arah kemiringan} rata-rata:
$2.94/\sqrt{12} = \ang{0.85}$.

\textbf{13.} Tidak: \ang{2.9} \omterm{def:b1:geological-maps:attitude}{arah kemiringan} pada perlapisan yang miring \ang{12} hanya
$2.9\sin\ang{12.1} = \ang{0.6}$ putaran bidang, lebih kecil daripada \ang{1.7} dari \omterm{def:b1:geological-maps:attitude}{kemiringan}. Ketidakrataan kecil pada permukaan perlapisan yang landai menyebarkan \omterm{def:b1:geological-maps:attitude}{arah kemiringannya} secara lebar.

\textbf{14.} $D = \ang{-2.5}$: $83.1 - 2.5 = \ang{080.6}$; rata-ratanya adalah
$080.6/12.1$, sama dengan milik ahli geologi pertama.

\textbf{15.} \omterm{def:b1:geological-maps:attitude}{Kemiringan} dibaca dengan klinometer, terhadap vertikal yang diberikan gravitasi; jarum tidak berperan.

\textbf{16.} Ya: putaran terhadap vertikal menggerakkan setiap pembacaan dan rata-ratanya dengan sudut yang sama (\cref{exo:b1:field-lab-methods-1:11}); jawabannya sama.

\textbf{17.} $10\times0.2 = \ang{2}$ galat pada setiap \omterm{def:b1:geological-maps:attitude}{arah kemiringan}, lebih dari dua kali galat baku \ang{0.85}: galat sistematis yang tidak dapat dideteksi oleh perataan.

\textbf{18.} $80.6 - 90 = \ang{350.6}$; orientasinya \omterm{def:b1:geological-maps:attitude}{jurus} \ang{351}, \omterm{def:b1:geological-maps:attitude}{kemiringan} \ang{12} ke timur.

\textbf{19.} \omterm{def:b1:geological-maps:attitude}{Arah kemiringan} $080.6$ dan $080.6$: tidak ada selisih pada ketelitian ini; \omterm{def:b1:geological-maps:attitude}{kemiringan} $12.1$ dan $12.2$: \ang{0.1}.

\textbf{20.} Normal $(\sin\ang{12.2}\sin\ang{80.6},
\sin\ang{12.2}\cos\ang{80.6}, \cos\ang{12.2})$ dan $\vec S/\|\vec S\|$ membentuk sudut \ang{0.13}.

\textbf{21.} $1.72/\sqrt{12} = \ang{0.50}$; selisihnya, \ang{0.13}, adalah seperempatnya: tidak bermakna. Kompas dan lubang bor sepakat.

\textbf{22.} Empat kali lebih banyak: 48 pembacaan.

\textbf{23.} Misalnya: pengaturan deklinasi yang keliru atau usang; kompas yang dibelokkan oleh besi di dekatnya (kendaraan, perkakas, batuan magnetik); klinometer yang tidak sesuai penyetelannya; pembacaan hanya pada permukaan yang paling tersingkap, yang mungkin tidak mewakili sayap itu. Perataan mengurangi galat acak, bukan galat sistematis.

\textbf{24.} \omterm{def:b1:field-lab-methods-1:mean-orientation}{Orientasi rata-rata} sayap timur adalah \textbf{081/12}: \omterm{def:b1:geological-maps:attitude}{arah kemiringan} \ang{081}, \omterm{def:b1:geological-maps:attitude}{kemiringan} \ang{12}.
\end{solution}
